% uid: ""
% school: "lycee"
% level: "1ere"
% chapter: "Activités numériques"
% title: "Un multiple de 30"
% difficulty: "Difficile"
% tags: ["divisibilite", "factorisation", "division-euclidienne", "chiffre-des-unites"]
% provenance: "own"
% batch: "2026-11-22-revision-1ere-A"

% ---STATEMENT---
Soit $n$ un entier naturel et $N=n^5-n$.
\begin{enumerate}
  \item Montrer que $N=(n-1)\,n\,(n+1)\,(n^2+1)$.
  \item Montrer que $N$ est divisible par $2$, puis par $3$, donc par $6$. [Distinguer les restes de la division de $n$ par $3$.]
  \item Montrer que $N$ est divisible par $5$. [Écrire $n=5k+r$ avec $0\leq r\leq4$ ; pour $r=2$ et $r=3$, regarder $n^2+1$.]
  \item En déduire que $N$ est divisible par $30$.
  \item Quel est le chiffre des unités de $2027^5$ ?
\end{enumerate}

% ---SOLUTION---
\begin{prop}[Divisibilité]
Si deux entiers \textbf{premiers entre eux} divisent $N$, leur produit divise $N$. Dans la division euclidienne de $n$ par $b$, \ib{$n=bk+r$ avec $0\leq r<b$}.
\end{prop}
\begin{enumerate}
  \item On factorise : $n^5-n=n\,(n^4-1)$. Puis $n^4-1=(n^2-1)(n^2+1)$ et $n^2-1=(n-1)(n+1)$ :
  \[N=\ib{$(n-1)\,n\,(n+1)\,(n^2+1)$}\]

  \item \textbf{Par $2$.} Parmi les deux entiers consécutifs $n-1$ et $n$, l'un est pair : $N$ est pair.

  \textbf{Par $3$.} On divise $n$ par $3$ : $n=3k$, $3k+1$ ou $3k+2$.

  \begin{itemize}
    \item $n=3k$ : le facteur $n$ est multiple de $3$.
    \item $n=3k+1$ : le facteur $n-1=3k$ est multiple de $3$.
    \item $n=3k+2$ : le facteur $n+1=3k+3=3(k+1)$ est multiple de $3$.
  \end{itemize}
  Dans les trois cas, $3\mid N$.

  $2$ et $3$ sont premiers entre eux : \ib{$6\mid N$}.

  \item On écrit $n=5k+r$ avec $r\in\{0,1,2,3,4\}$.

  \begin{itemize}
    \item $r=0$ : le facteur $n=5k$ est multiple de $5$.
    \item $r=1$ : le facteur $n-1=5k$ est multiple de $5$.
    \item $r=4$ : le facteur $n+1=5k+5=5(k+1)$ est multiple de $5$.
    \item $r=2$ : $n^2+1=25k^2+20k+4+1=5\,(5k^2+4k+1)$ est multiple de $5$.
    \item $r=3$ : $n^2+1=25k^2+30k+9+1=5\,(5k^2+6k+2)$ est multiple de $5$.
  \end{itemize}
  Dans tous les cas, \ib{$5\mid N$}.

  \item $5$ et $6$ sont premiers entre eux, et $6\mid N$, $5\mid N$ : $N$ est divisible par $5\times6=\boxed{30}$.

  \item Comme $30\mid N$, on a aussi $10\mid N$ : $n^5-n=10m$, donc
  \[n^5=n+10m\]
  Les nombres $n^5$ et $n$ ont le \textbf{même chiffre des unités}.

  Pour $n=2027$ : le chiffre des unités de $2027^5$ est \ib{$7$}.
\end{enumerate}
\textbf{Idée :} factoriser, puis traiter chaque diviseur premier ($2$, $3$, $5$) par les restes de la division euclidienne.

\textbf{Piège :} conclure \og divisible par $30$ \fg{} à partir de $2$, $3$, $5$ sans passer par des entiers premiers entre eux (par exemple $4$ et $6$ divisent $12$, mais pas $24$).
